\chapter{System Design}

This chapter presents the architecture of \updesk{}. It begins with the overall
component decomposition, then treats the signaling protocol, the
connection-establishment sequence, the data model and the security reasoning in
turn.

\section{Architectural Overview}

\updesk{} is composed of three cooperating roles and one coordinating service.
The \emph{controller} is the machine from which a user drives a remote session.
The \emph{host} is the machine being controlled; it exists in two forms, a
graphical Tauri application for attended use and a headless native binary for
unattended use. The \emph{signaling server} is a lightweight Rust process that
introduces controllers to hosts and relays the messages needed to negotiate a
peer connection, but carries no media. \figref{fig:arch} shows the arrangement.

\begin{figure}[H]
\centering
\begin{tikzpicture}[
  node distance=2.4cm and 3.2cm,
  box/.style={draw,rounded corners,minimum width=3cm,minimum height=1.1cm,align=center,font=\small},
  server/.style={draw,rounded corners,fill=black!8,minimum width=3.2cm,minimum height=1.1cm,align=center,font=\small},
  sig/.style={-{Latex[length=2mm]},dashed,thick},
  media/.style={-{Latex[length=2.6mm]},very thick}
]
  \node[server] (srv) {Signaling Server\\{\scriptsize (Rust, WSS)}};
  \node[box, below left=of srv] (ctl) {Controller\\{\scriptsize (Tauri app)}};
  \node[box, below right=of srv] (host) {Host\\{\scriptsize (native / Tauri)}};

  \draw[sig] (ctl) -- (srv) node[midway,above,sloped,font=\scriptsize] {auth + SDP/ICE};
  \draw[sig] (host) -- (srv) node[midway,above,sloped,font=\scriptsize] {auth + register};
  \draw[media] (host) -- (ctl)
      node[midway,below,font=\scriptsize,align=center] {direct WebRTC\\DTLS-SRTP media + control};

  \node[font=\scriptsize,align=center,below=1.4cm of srv] {\textit{server sees metadata only ---}\\\textit{never the media keys}};
\end{tikzpicture}
\caption{Component architecture. Dashed lines are signaling; the solid line is
the direct, end-to-end encrypted peer connection that carries screen frames and
control events.}
\label{fig:arch}
\end{figure}

The essential property visible in the figure is that the two solid endpoints
exchange media directly. The signaling server sits on the dashed paths only. It
authenticates both endpoints and passes their session descriptions back and
forth, but once the peer connection is live the server can be removed without
disturbing the session.

\section{The Native Host Internals}

The native host is the most involved component because it must do the work a
browser would otherwise do, without a browser. \figref{fig:nativehost} shows its
internal pipeline. A capture thread pulls framebuffers from the Desktop
Duplication API; each frame is handed to an H.264 encoder; encoded access units
are written onto a WebRTC video track; and, in the reverse direction, control
messages arriving on a data channel are decoded into synthetic input events and
injected into the desktop.

\begin{figure}[H]
\centering
\begin{tikzpicture}[
  stage/.style={draw,rounded corners,minimum width=2.5cm,minimum height=1cm,align=center,font=\scriptsize},
  flow/.style={-{Latex[length=2mm]},thick},
  node distance=0.9cm and 1.2cm
]
  \node[stage] (cap) {Capture\\{\tiny scrap / DXGI}};
  \node[stage,right=of cap] (enc) {Encode\\{\tiny OpenH264}};
  \node[stage,right=of enc] (track) {Video Track\\{\tiny webrtc-rs}};
  \node[stage,right=of track] (pc) {Peer\\Connection};

  \node[stage,below=1.6cm of track] (inj) {Input Inject\\{\tiny SendInput}};
  \node[stage,left=of inj] (dc) {Control\\Channel};

  \draw[flow] (cap) -- (enc);
  \draw[flow] (enc) -- (track);
  \draw[flow] (track) -- (pc);
  \draw[flow] (pc.south) |- (dc);
  \draw[flow] (dc) -- (inj);

  \node[draw,dotted,fit=(cap)(enc)(track)(pc)(dc)(inj),inner sep=10pt,label={[font=\scriptsize]above:\textit{native-host process}}] {};
\end{tikzpicture}
\caption{Native host pipeline. The upper row is the outbound video path; the
lower row is the inbound control path.}
\label{fig:nativehost}
\end{figure}

\section{Signaling Protocol}

All signaling messages are JSON objects distinguished by a \texttt{type} field
and carried over a WebSocket (plain \texttt{ws://} in development,
\texttt{wss://} in deployment). \tabref{tab:messages} enumerates the message
types actually implemented by the server and the native host.

\begin{table}[H]
\centering
\caption{Signaling message types}
\label{tab:messages}
\small
\begin{tabular}{@{}p{3.2cm}p{2cm}p{7.6cm}@{}}
\toprule
\textbf{Message} & \textbf{Direction} & \textbf{Purpose} \\
\midrule
\texttt{auth\_init}      & server$\to$peer & Server issues a random challenge nonce. \\
\texttt{auth\_response}  & peer$\to$server & Peer returns its public key and an Ed25519 signature over the nonce. \\
\texttt{register}        & host$\to$server & Host announces availability under its connect identifier. \\
\texttt{connect\_request}& controller$\to$server & Controller asks to reach a host by connect identifier. \\
\texttt{session\_response}& host$\to$server & Host accepts or rejects the incoming request. \\
\texttt{offer}           & host$\to$controller & SDP offer, relayed by the server. \\
\texttt{answer}          & controller$\to$host & SDP answer, relayed by the server. \\
\texttt{ice\_candidate}  & both & An ICE candidate, relayed as it is discovered. \\
\texttt{end\_session}    & both & Orderly teardown of the signaling association. \\
\bottomrule
\end{tabular}
\end{table}

\subsection{Enrollment}

Before a machine can authenticate it must enroll. Enrollment is gated by a
single-use code of the form \texttt{DEV0-TEST:device} or
\texttt{CTL0-TEST:controller}, seeded in the server's configuration and consumed
on first use. During enrollment the peer generates an Ed25519 key pair, keeps
the private key locally, and registers the public key against the code. From
that point the public key \emph{is} the machine's identity; the code is spent
and cannot be reused. This deliberately prevents anonymous participation: an
attacker cannot register a host or controller without first obtaining a valid
code out of band.

\subsection{Authentication Handshake}

Every WebSocket connection begins with the same challenge--response exchange,
shown in \figref{fig:auth}, before any other message is honoured. The server
sends \texttt{auth\_init} carrying a fresh random nonce. The peer signs that
nonce with its Ed25519 private key and replies with \texttt{auth\_response}
containing its public key and the signature. The server verifies the signature
against the enrolled key. Because the nonce is random and single-use, a captured
\texttt{auth\_response} cannot be replayed against a later challenge.

\begin{figure}[H]
\centering
\begin{tikzpicture}[
  font=\scriptsize,
  actor/.style={draw,minimum width=2.4cm,minimum height=0.8cm},
  msg/.style={-{Latex[length=2mm]}}
]
  \node[actor] (p) {Peer};
  \node[actor,right=5.5cm of p] (s) {Server};
  \draw[thick] (p.south) -- ++(0,-4.6);
  \draw[thick] (s.south) -- ++(0,-4.6);
  \coordinate (p0) at ($(p.south)+(0,-0.5)$);
  \coordinate (s0) at ($(s.south)+(0,-0.5)$);
  \draw[msg] ($(s0)$) -- ($(p0)$) node[midway,above]{\texttt{auth\_init}\{ nonce \}};
  \coordinate (p1) at ($(p.south)+(0,-1.6)$);
  \node[right=0.1cm of p1,align=left,font=\tiny] {sign(nonce, sk)};
  \coordinate (p2) at ($(p.south)+(0,-2.4)$);
  \coordinate (s2) at ($(s.south)+(0,-2.4)$);
  \draw[msg] (p2) -- (s2) node[midway,above]{\texttt{auth\_response}\{ pk, sig \}};
  \coordinate (s3) at ($(s.south)+(0,-3.2)$);
  \node[left=0.1cm of s3,align=right,font=\tiny] {verify(sig, pk, nonce)\\lookup enrolled pk};
  \coordinate (p4) at ($(p.south)+(0,-4.1)$);
  \coordinate (s4) at ($(s.south)+(0,-4.1)$);
  \draw[msg] (s4) -- (p4) node[midway,above]{authenticated};
\end{tikzpicture}
\caption{Ed25519 challenge--response authentication performed on every
connection.}
\label{fig:auth}
\end{figure}

\section{Connection Establishment}

Once both endpoints are authenticated, establishing a session follows the
sequence in \figref{fig:session}. The host registers and waits. A controller
issues a \texttt{connect\_request} naming the host's nine-digit connect
identifier. The server forwards the request; the host answers with
\texttt{session\_response}. On acceptance the WebRTC negotiation begins: the host
emits an SDP \texttt{offer}, the controller replies with an \texttt{answer}, and
both sides trickle \texttt{ice\_candidate} messages as ICE discovers transport
addresses. When a candidate pair succeeds, the DTLS handshake runs directly
between the peers, SRTP keys are derived, and media begins to flow on the direct
path. The server has by then done its entire job.

\begin{figure}[H]
\centering
\begin{tikzpicture}[
  font=\scriptsize,
  actor/.style={draw,minimum width=2.2cm,minimum height=0.8cm},
  msg/.style={-{Latex[length=2mm]}},
  self/.style={-{Latex[length=2mm]}}
]
  \node[actor] (c) {Controller};
  \node[actor,right=3.4cm of c] (s) {Server};
  \node[actor,right=3.4cm of s] (h) {Host};
  \foreach \n in {c,s,h}{\draw[thick] (\n.south) -- ++(0,-6.4);}

  \newcommand{\line}[4]{%
    \coordinate (a) at ($(#1.south)+(0,-#3)$);
    \coordinate (b) at ($(#2.south)+(0,-#3)$);
    \draw[msg] (a) -- (b) node[midway,above]{#4};}

  \line{h}{s}{0.6}{\texttt{register}}
  \line{c}{s}{1.4}{\texttt{connect\_request}(id)}
  \line{s}{h}{2.1}{forward}
  \line{h}{s}{2.9}{\texttt{session\_response}(accept)}
  \line{h}{s}{3.7}{\texttt{offer} (SDP)}
  \line{s}{c}{3.7}{}
  \line{c}{s}{4.5}{\texttt{answer} (SDP)}
  \line{s}{h}{4.5}{}
  \line{c}{h}{5.3}{\texttt{ice\_candidate} $\leftrightarrow$ (via server)}
  \draw[msg,very thick] ($(h.south)+(0,-6.1)$) -- ($(c.south)+(0,-6.1)$)
      node[midway,above]{\textbf{direct DTLS-SRTP media}};
\end{tikzpicture}
\caption{Session-establishment sequence. After the final step the media path is
direct and the server is no longer involved.}
\label{fig:session}
\end{figure}

\section{Security Model}

The security posture of \updesk{} rests on a clean division of responsibility.

\textbf{Identity is a key, not a password.} Each participant is named by an
Ed25519 public key whose private half never leaves the machine. Compromising the
server does not yield any participant's signing key.

\textbf{Enrollment is closed.} A public key is accepted only in exchange for a
single-use code, so the population of legitimate endpoints is controlled and
auditable.

\textbf{Freshness defeats replay.} The per-connection random nonce ensures that
recording a valid authentication exchange confers no ability to authenticate
later.

\textbf{The server cannot read the session.} DTLS-SRTP keys are negotiated
directly between the peers. The server relays the SDP that begins that
negotiation but never possesses the derived keys, so it is structurally unable
to decrypt frames or control events. This is the property that distinguishes the
design from relay-mediated commercial products.

\textbf{Unattended acceptance is host-verified.} Automatic acceptance of an
unattended session is conditioned on a password checked by the host itself, not
by the server, so the server never sees the unattended secret either.

\tabref{tab:threats} maps the principal threats to the mechanisms that address
them.

\begin{table}[H]
\centering
\caption{Threat model summary}
\label{tab:threats}
\small
\begin{tabular}{@{}p{5.4cm}p{8cm}@{}}
\toprule
\textbf{Threat} & \textbf{Mitigation} \\
\midrule
Unauthorised host or controller joins & Closed enrollment via single-use codes; Ed25519 identity. \\
Impersonation of an enrolled peer & Signature over a server-issued nonce, verified against the enrolled key. \\
Replay of a captured handshake & Fresh random nonce per connection. \\
Eavesdropping by the server or a relay & End-to-end DTLS-SRTP; keys never held by the server. \\
Passive network eavesdropping & \texttt{wss://} for signaling; SRTP for media. \\
Abuse of unattended acceptance & Host-verified unattended password. \\
\bottomrule
\end{tabular}
\end{table}
